\scanpage{12}
Indeed, if $Y$ is a tame space, a \deleted{bijection} \unsure{arrow} from the inductive system $(X_i)$ (in $\mathcal E$) to $Y$, \inserted{$(f_i:X_i\mathbin{-}Y)_{i\in I}$}, is equivalent to giving a map $f:X\mathbin{-}Y$, such that the $f\circ g_i$ are tame. It is \unsure{obvious} : let us prove that this implies that $f$ is tame, i.e. that its graph $\Gamma_f\subset X\times Y$ is a tame subset. Now $X$ being the union of the images of the $X_i$, it is the case that $\Gamma_f$ is the union of the images of the $\Gamma_{f_i}\subset X_i\times Y$ under
\[
 X_i\unsure{\deleted{\times Y}|\times Y}\xrightarrow{g_i\times\operatorname{id}_Y}X\times Y.
\]
But these images are tame, and therefore their union too.

\uline{Corollary} Let $f:X\longrightarrow Y$ be a tame map, $\mathcal R_f$ the equivalence relation defined by $f$ (which is therefore a tame subset of $X\times X$, and is an $\mathcal E$-equivalence relation). Then $X/\mathcal R_f$ exists in $\mathcal E=\operatorname{Esp}_{(\mathcal M_*)}$, and is \inserted{canonically} isomorphic to $f(X)$ (with induced structure)

\deleted{\uline{Corollary 2} Consider}
\deleted{\diagram{\begin{tikzcd}[column sep=small]
&\deleted{Z}\arrow[dl,hook',"\deleted{i_1}"']\arrow[dr,hook,"\deleted{i_2}"]&\\
\deleted{X_1}&&\deleted{X_2}
\end{tikzcd}}}
\deleted{in $\mathcal E$.}
\mspage{11}
